\documentclass[11pt,a4paper]{article}

\usepackage[margin=1in]{geometry}
\usepackage{amsmath,amssymb,amsthm,mathtools}
\usepackage{algorithm}
\usepackage{algpseudocode}
\usepackage{booktabs}
\usepackage{graphicx}
\usepackage{hyperref}
\usepackage{cleveref}
\usepackage{enumitem}
\usepackage{microtype}

\graphicspath{{../figures/}}

\newtheorem{theorem}{Theorem}[section]
\newtheorem{lemma}[theorem]{Lemma}
\newtheorem{proposition}[theorem]{Proposition}
\newtheorem{corollary}[theorem]{Corollary}
\newtheorem{definition}[theorem]{Definition}
\newtheorem{remark}[theorem]{Remark}
\newtheorem{construction}[theorem]{Construction}

\newcommand{\Enc}{\mathsf{Encode}}
\newcommand{\Det}{\mathsf{Detect}}
\newcommand{\KeyGen}{\mathsf{KeyGen}}
\newcommand{\Embed}{\mathsf{Embed}}
\newcommand{\Extract}{\mathsf{Extract}}
\newcommand{\negl}{\mathsf{negl}}
\newcommand{\wt}{\mathrm{wt}}
\newcommand{\Bern}{\mathrm{Bern}}
\newcommand{\BSC}{\mathrm{BSC}}
\newcommand{\PRCC}{\ensuremath{\mathrm{PRCC}}}
\newcommand{\sk}{\mathit{sk}}

\title{Pseudorandom Cyclic Codes:\\
Syntax, a Block-Circulant Candidate, and a Lattice Watermark}
\author{
Anonymous / Working Draft\\
\texttt{PRCC project}
}
\date{\today}

\begin{document}
\maketitle

\begin{abstract}
We study zero-bit \emph{Pseudorandom Cyclic Codes} (\PRCC{}): a keyed coding syntax
for the simultaneous ring-shift \(T_k\), with desired completeness, uniform-string
soundness, single-sample pseudorandomness, and \(\BSC_p\) robustness.
A block-circulant candidate encodes the coset \(Hc=\mathbf{1}\) and detects by
syndrome weight.
The identity \(HT_k=R_kH\) is the standard quasi-cyclic equivariance; it makes
integer sector rolls free at the detector.
We then instantiate that candidate as a post-hoc image watermark: two-coset
quantization index modulation (QIM) on hard radial mean-difference coefficients,
plus a search over two content centers and \(L\) residual angles.
That instance is an application, not a proof that the candidate meets every
\PRCC{} axiom.

At a tested operating point (COCO images bilinearly resized to \(256\times256\),
median PSNR \(\approx 38\,\mathrm{dB}\), encode noise \(\eta=0\), two centers
\(\times L=8\)), five keys on a hold-out split meet TPR \(\ge 90\%\) under random
rotation with JPEG quality~70, with Clopper--Pearson 95\% FPR upper bounds
\(\le 1\%\) (Q70-calibrated thresholds).
Relative to a matched PRC detector that ORs over integer shifts, the construction
cuts code-detection calls by \(64\times\) (16 vs.\ 1024 per image).
Against a keyed \(\pm 1\) template with FFT cyclic correlation on the same carrier and
geometry, detect-call counts match (16 vs.\ 16) and the correlator attains higher TPR;
the measured structural win is therefore against \(m\)-fold shift-OR code search, not
against FFT-amortized template correlation.
Block-circulant equivariance and QIM are tools.
The coding contribution is the \PRCC{} syntax together with a candidate whose
integer-shift closure is proved and whose pseudorandomness is not.
The watermark measurements are application metrics at \(\eta=0\); that operating
point has a public parity constraint and is not a pseudorandom \PRCC{} realization.
We do not claim image-level undetectability.
\end{abstract}

\section{Introduction}
\label{sec:intro}

Rotation remains awkward for zero-bit image watermarks.
Classical geometric methods often map rotation to a cyclic shift of an extracted
1-D descriptor and then search or correlate~\cite{lin2001rst}.
Fourier-ring schemes for diffusion images similarly rely on angular
structure~\cite{wen2023treering,ci2024ringid}, while cryptographic pseudorandom
codes (PRCs)~\cite{christ2024prc,gunn2024undetectable} give strong substitution
guarantees but treat cyclic shift as an external OR over \(m\) detector calls.

The paper is organized in three layers.
First we fix the \PRCC{} syntax and the four target properties (completeness,
uniform-string soundness, pseudorandomness, BSC robustness).
Second we give a block-circulant candidate and state, property by property, what
is proved and what is not.
Third we embed that candidate in a spatial lattice watermark and report distortion,
rotation+JPEG detection, and detection-call counts.
The algebraic identity \(HT_k=R_kH\) is a standard quasi-cyclic
fact~\cite{townsend1967cyclic,fossorier2004qc}; two-coset embedding is
QIM~\cite{chen2001qim}.
Refusing to claim that the \(\eta=0\) instance is pseudorandom does not remove the
definition: that instance is a candidate boundary, not a substitute for the syntax.

\paragraph{Contributions.}
\begin{itemize}[leftmargin=*]
  \item A zero-bit \PRCC{} definition under simultaneous ring shifts \(T_k\),
        including the ring-OR PRC baseline that any useful construction should beat
        on detection cost or end-to-end metrics.
  \item A block-circulant candidate (\(Hc=\mathbf{1}\), syndrome detection) with an
        equivariance proof, a uniform-string soundness bound under full rank, and an
        explicit gap: \(\eta=0\) fails pseudorandomness by a public parity constraint.
  \item An image instantiation at \(\eta=0\), \(w=1\): hard radial mean-difference
        QIM, two-center residual search, and COCO measurements (median PSNR
        \(\approx 38\,\mathrm{dB}\), five-key hold-out confirmation, \(64\times\)
        fewer code calls than shift-OR PRC, and a stronger template+FFT correlator
        on the same geometry).
\end{itemize}

\paragraph{Non-claims.}
Image TPR/FPR numbers are application metrics.
They are not proofs of uniform-string soundness, BSC robustness, or
pseudorandomness.
A failed full-circulant square matrix is deferred to~\Cref{app:failed}.

\section{Related Work}
\label{sec:related}

\paragraph{Classical geometric watermarking.}
Lin et~al.~\cite{lin2001rst} build an RST-resilient watermark from a
Fourier--Mellin-style 1-D projection: image rotation induces a cyclic shift of the
extracted signal, detected by search and correlation, with reported false positives,
rotation, and JPEG tests.
Our setting is likewise post-hoc zero-bit watermarking under rotation and JPEG, so
that line of work is the closest classical reference.
We differ in payload geometry (concentric spatial annuli rather than log-polar
Fourier magnitudes) and in replacing template correlation by a shift-invariant
syndrome test at the coding layer.

\paragraph{QIM and quasi-cyclic codes.}
Chen and Wornell~\cite{chen2001qim} introduce quantization index modulation; our
carrier uses a standard two-coset dithered quantizer on scalar coefficients.
Quasi-cyclic / block-circulant parity-check structure is classical in coding
theory~\cite{townsend1967cyclic,fossorier2004qc,misoczki2013mdpc}.
We use \(HT_k=R_kH\) as an off-the-shelf equivariance tool, not as a new code-family
claim.

\paragraph{PRCs and generative watermarks.}
Christ and Gunn~\cite{christ2024prc} define PRCs; variants
follow~\cite{alrabiah2025ideal,christ2025permuted,golowich2024edit}.
Image PRC watermarks embed in latents~\cite{gunn2024undetectable}.
A ring-OR PRC baseline answers how much the coding layer saves versus searching
integer shifts; it does not by itself locate the method among classical geometric
watermarks.

\paragraph{Diffusion watermarks and rotation.}
Tree-Ring~\cite{wen2023treering} and RingID~\cite{ci2024ringid} analyze Fourier-ring
embedding in generative pipelines.
Fractional shifts and cropping can destroy intended robustness~\cite{ci2024ringid}.
Our experiments are post-hoc on natural images with a spatial carrier, not
claims about diffusion inversion.

\section{Preliminaries}
\label{sec:prelim}

Fix \(b\ge 2\) concentric rings and \(m\) angular samples per ring.
Codewords live in \(\{0,1\}^n\) with \(n=bm\).
The length-\(m\) cyclic shift is \((R_k x)_i=x_{i-k\bmod m}\).
The simultaneous ring action is
\[
T_k=\operatorname{diag}(R_k,\ldots,R_k)\in\{0,1\}^{n\times n},\qquad k\in\mathbb{Z}/m\mathbb{Z}.
\]
Write \(\wt(\cdot)\) for Hamming weight.
\(\BSC_p\) flips each bit independently with probability \(p\).
One angular bin has width \(\Delta=360^\circ/m\).
An integer sector roll corresponds to some \(T_k\); a residual angle
\(\delta\in[0,\Delta)\) accounts for misalignment within a bin.
Worst-case \(p\)-bounded adversarial flips are deferred.

\section{PRCC: Syntax and Desired Properties}
\label{sec:prcc}

\begin{definition}[Zero-bit \PRCC{} under \(T_k\)]
\label{def:prcc}
A triple \((\KeyGen,\Enc,\Det)\) is a \emph{zero-bit Pseudorandom Cyclic Code} if
there exists \(p\in(0,1/2)\) such that the following hold.
\(\KeyGen(1^\lambda)\) outputs a secret key \(\sk\).
\(\Enc(\sk)\) outputs a codeword in \(\{0,1\}^n\).
\(\Det(\sk,y)\) outputs a bit.

\paragraph{(C) Completeness / detect-sense \(T_k\)-closure.}
For all \(k\in\mathbb{Z}/m\mathbb{Z}\),
\[
\Pr\bigl[\Det(\sk, T_k(\Enc(\sk)))=1\bigr]\ge 1-\negl(\lambda).
\]

\paragraph{(S) Soundness against uniform strings.}
\[
\Pr\bigl[\Det(\sk,U)=1\bigr]\le\negl(\lambda)
\quad\text{for }U\leftarrow\{0,1\}^n.
\]
This does not imply a low false-positive rate on bits extracted from unmarked
images, which need not be uniform.
A homogeneous syndrome test accepts the all-zero string.

\paragraph{(P) Pseudorandomness (single-sample).}
For every PPT adversary \(\mathcal{A}\) that does not receive \(\sk\),
\[
\bigl|\Pr[\mathcal{A}(\Enc(\sk))=1]-\Pr[\mathcal{A}(U)=1]\bigr|\le\negl(\lambda).
\]
Multi-sample pseudorandomness under one \(\sk\) is a comparison target, not part of
this definition~\cite{christ2024prc}.
Bit pseudorandomness does not imply image-level undetectability.

\paragraph{(R) Robustness under \(\BSC_p\).}
\[
\forall k\in\mathbb{Z}/m\mathbb{Z}:\quad
\Pr\bigl[\Det(\sk,\BSC_p(T_k(\Enc(\sk))))=1\bigr]\ge 1-\negl(\lambda).
\]
\end{definition}

\begin{remark}[Trivial baseline]
\label{rem:ringor}
Given any zero-bit PRC detector \(D_{\mathrm{PRC}}\) for length \(n=bm\), define
\[
D_{\mathrm{ring}}(y)=\bigvee_{k=0}^{m-1} D_{\mathrm{PRC}}(T_k y).
\]
The encoder is unchanged, so (P) is inherited.
(C) holds for exact \(T_k\).
Under \(\BSC_p\), (R) inherits up to the original completeness gap.
If the base false-positive rate on uniform strings is \(\varepsilon\), the ring-OR
rate is at most \(m\varepsilon\).
A useful \PRCC{} must beat this baseline on detection cost or on end-to-end
watermark metrics under matched false-positive rate and attack strength.
\end{remark}

\section{Candidate: Block-Circulant Syndromes}
\label{sec:method}

\subsection{Key generation and coset encoding}
\label{sec:code}

\begin{construction}[Block-circulant zero-bit code]
\label{cons:block}
\textbf{Parameters:} rings \(b\), angular length \(m\), common block row weight \(w\),
encode noise \(\eta\), threshold \(\tau\).

\paragraph{$\KeyGen$.}
For \(j=1,\ldots,b\), sample a generator \(h^{(j)}\in\{0,1\}^m\) of weight \(w\).
Let \(A_j\) be the \(m\times m\) circulant with first row \(h^{(j)}\), and set
\(H=[A_1\mid\cdots\mid A_b]\in\{0,1\}^{m\times n}\).
Reject and resample until \(A_b\) is invertible over the circulant algebra
(equivalently \(\gcd(h^{(b)}(X),X^m-1)=1\)).
Output \(\sk=(H,\tau,\eta)\).

\paragraph{$\Enc$.}
Sample free blocks \(c^{(1)},\ldots,c^{(b-1)}\leftarrow\{0,1\}^m\) uniformly.
Solve \(A_b c^{(b)}=\mathbf{1}+\sum_{j<b}A_j c^{(j)}\) for \(c^{(b)}\).
Output \(c=(c^{(1)}\|\cdots\|c^{(b)})\oplus e\) with \(e\leftarrow\Bern(\eta)^n\).
At the performance operating point we take \(\eta=0\).

\paragraph{Syndrome detection.}
Define \(s(y)=\wt(Hy+\mathbf{1})\).
The bitstring detector accepts iff \(s(y)\le\tau\).
The near-constant block rule of~\Cref{sec:const} is part of the image detector
only; it is not part of this bitstring syntax.
\end{construction}

\begin{proposition}[Integer-shift equivariance]
\label{prop:inv}
Each circulant block satisfies \(A_j R_k=R_k A_j\), hence \(HT_k=R_kH\).
If \(Hc=\mathbf{1}\) then \(H(T_kc)=R_k\mathbf{1}=\mathbf{1}\), so
\(s(T_kc)=0\).
\end{proposition}

\begin{proposition}[Kernel dimension]
\label{prop:ker}
Always \(\dim\ker H\ge (b-1)m\).
If \(A_b\) is invertible then \(\mathrm{rank}(H)=m\) and \(\dim\ker H=(b-1)m\).
\end{proposition}

\begin{proposition}[Uniform-string soundness under full rank]
\label{prop:S}
If \(\mathrm{rank}(H)=m\) and \(U\leftarrow\{0,1\}^n\), then \(HU\) is uniform on
\(\mathbb{F}_2^m\).
Thus \(\wt(HU+\mathbf{1})\sim\mathrm{Bin}(m,1/2)\) and
\[
\Pr[\Det(\sk,U)=1]=\Pr[\mathrm{Bin}(m,1/2)\le\tau].
\]
Choosing \(\tau\ll m/2\) makes this probability exponentially small in \(m\).
This is property~(S) for the bitstring detector.
It is not a statement about unmarked images.
\end{proposition}

\paragraph{The \(w=1\) operating point.}
When \(w=1\), each generator is a standard basis vector, so every \(A_j\) is a
circulant permutation matrix (a pure cyclic shift).
Writing ring blocks \(y^{(1)},\ldots,y^{(b)}\in\{0,1\}^m\) and secret offsets
\(t_j\) for the nonzero of \(h^{(j)}\),
\[
(Hy)_i=\bigoplus_{j=1}^{b} y^{(j)}_{i-t_j\bmod m}.
\]
Thus \(s(y)\) counts, after aligning each ring by its secret shift, how many
angular positions violate the cross-ring parity that targets the all-ones syndrome.
Detection is a shift-aligned parity-check count, not iterative decoding.
The weight \(w=1\) is the application operating point used below; it is not part of
\Cref{def:prcc}.

\subsection{What the candidate establishes}
\label{sec:status}

\paragraph{Established.}
\Cref{prop:inv} gives \(HT_k=R_kH\).
If \(\eta=0\) and \(Hc=\mathbf{1}\), then \(s(T_kc)=0\), so completeness~(C) holds
with probability \(1\) for exact integer shifts of noiseless codewords.
\Cref{prop:S} gives the uniform-string bound for~(S) whenever \(A_b\) is invertible
and \(\tau\) is set below the binomial tail.

\paragraph{Not established.}
When \(\eta=0\), summing coordinates yields
\(\sum_j h^{(j)}(1)\,c^{(j)}(1)=m\bmod 2\).
Odd-weight generators, including every \(w=1\) generator, have \(h^{(j)}(1)=1\), so
the block parities of a noiseless codeword are constrained, while a uniform string
meets the same relation with probability \(1/2\).
The \(\eta=0\) candidate therefore fails~(P) information-theoretically.
The image experiments do not prove~(R) at a \(\negl(\lambda)\) level, and a low
false-positive rate on extracted image bits does not replace~(S).
Pixel rotation is not proved equal to \(T_k\) plus light noise; that correspondence
is measured only in the application below.

\section{Application: Spatial Lattice Watermark}
\label{sec:app}

The sections below instantiate~\Cref{cons:block} on images.
They measure embedding distortion, detection rate under rotation and JPEG, and
detection-call counts against a shift-OR PRC and a template correlator.
Those are application metrics.
They are not proofs of~(C), (S), (P), or~(R).
A stronger correlation baseline limits the competitiveness of this watermark
instance; it is not a reason to drop the \PRCC{} syntax.
Images are COCO photographs bilinearly resized to \(256\times256\) RGB.
Rotation attacks use pad-and-crop, optionally followed by JPEG.
The reported operating point uses \(\eta=0\).

\subsection{Near-constant block rejection}
\label{sec:const}

If every length-\(m\) block of \(y\) has mean bit fraction in
\([0,\rho]\cup[1-\rho,1]\) with \(\rho=0.05\), the detector returns ``absent''
before comparing \(s(y)\) to \(\tau\).
The rule is \(T_k\)-invariant and enabled in all reported image experiments.
Reproducing reported FPR/TPR requires it.

\subsection{Lattice carrier (two-coset QIM)}
\label{sec:lattice}

\paragraph{Coefficients.}
Each annulus \([r_0,r_1]\) is split at the midpoint \(r_{\mathrm{mid}}=(r_0+r_1)/2\)
into hard inner and outer radial bands (no soft radial weights).
Angular sector \(t\) keeps only pixels whose fractional sector coordinate lies in
\([\tfrac{1}{2}g,\,1-\tfrac{1}{2}g)\) with angular guard \(g=0.15\), excluding
sector boundaries.
Let \(\mu^{\mathrm{out}}_{j,t}\) and \(\mu^{\mathrm{in}}_{j,t}\) be the mean gray
values (unit scale) in the outer and inner bands of ring \(j\), sector \(t\).
The scalar coefficient is the mean difference
\[
v_{j,t}=\mu^{\mathrm{out}}_{j,t}-\mu^{\mathrm{in}}_{j,t}
\]
when both bands are nonempty, else \(0\).

\paragraph{Two-coset QIM~\cite{chen2001qim}.}
With lattice spacing \(s>0\), bit \(b\in\{0,1\}\) selects the coset
\(\{(2k+b)\,s/2:k\in\mathbb{Z}\}\).
Embedding replaces \(v\) by the nearest point \(\widehat{v}\) of that coset:
\[
\widehat{v}
=\operatorname{round}\!\bigl((v-b\cdot s/2)/s\bigr)\,s+b\cdot s/2.
\]
Let \(\delta=\widehat{v}-v\).
Pixels in the inner (resp.\ outer) band of sector \(t\) receive additive gray
perturbations
\[
d^{\mathrm{in}}_t=-\frac{n^{\mathrm{out}}_t}{n^{\mathrm{in}}_t+n^{\mathrm{out}}_t}\,\delta_t,
\qquad
d^{\mathrm{out}}_t=\frac{n^{\mathrm{in}}_t}{n^{\mathrm{in}}_t+n^{\mathrm{out}}_t}\,\delta_t,
\]
where \(n^{\mathrm{in/out}}_t\) are pixel counts; this realizes the target change in
\(v\) with minimum squared pixel change among constant-per-band pushes.
RGB images apply the same gray residual to each channel, then clip to the legal
range.
Extraction returns the closer coset as a hard bit.

The locked operating point uses \(m=64\), \(b=4\), \(w=1\), \(\eta=0\), and
\(s\approx 0.1416\) (highest median-PSNR setting that passed the gates
of~\Cref{sec:protocol} on a development key;~\Cref{tab:psnr-scan}).

\subsection{Detection: centers, residual angles, and decision}
\label{sec:detect}

\begin{algorithm}[t]
\caption{Blind detection (baseline~C)}
\label{alg:detect}
\begin{algorithmic}[1]
\Require image \(I\), key \(\sk\), layout, lattice step \(s\), centers
\(\mathcal{C}=\{(127.5,127.5),(128,128)\}\), residual grid size \(L\)
\State \(\Delta\leftarrow 360^\circ/m\); \(\mathcal{D}\leftarrow\{i\Delta/L:i=0,\ldots,L-1\}\)
\State \(s^\star\leftarrow +\infty\)
\For{each center \(c\in\mathcal{C}\)}
  \For{each \(\delta\in\mathcal{D}\)}
    \State \(y\leftarrow\Extract(I;\,c,\delta,s)\) \Comment{hard bits, length \(n\)}
    \If{\(y\) fails near-constant reject}
      \State \textbf{continue}
    \EndIf
    \State \(s^\star\leftarrow\min(s^\star,\,s(y))\)
  \EndFor
\EndFor
\State \Return \(\mathbf{1}[s^\star\le\tau]\)
\end{algorithmic}
\end{algorithm}

Pad-and-crop rotation on a \(256\times256\) canvas expands to an even or odd side
length depending on angle.
Even expansion preserves content center \((127.5,127.5)\); odd expansion moves it to
\((128,128)\).
Searching both centers removes the half-pixel residual-angle valley of
default-center extraction~(\Cref{tab:center-sync}).
This is not general translation synchronization.

With \(L=8\), each image incurs \(2\times L=16\) extractions and 16 syndrome
evaluations; the integer sector index never appears in the search.

\subsection{Threshold calibration}
\label{sec:calib}

Every key freezes a single threshold before test evaluation.

\paragraph{Q70-only rule.}
Calibrate on unmarked scores and on random-rotation+JPEG~70 attacked scores
(empirical FPR strictly below \(0.5\%\) on both, tune-set TPR \(\ge 90\%\), largest
feasible integer \(\tau\) for PRCC).
This is the rule used by the hold-out confirmation table, the PRC comparison, and
the template+correlation comparison.

\paragraph{Mixed-Q rule (formal multi-quality protocol).}
On calibration, aggregate attacked negatives \emph{per image} by the minimum score
over \(Q\in\{50,60,70,80,90\}\) (worst case), and tune positives by the maximum
score over \(Q\).
Freeze one \(\tau\) and evaluate every \(Q\) without retuning.
Implementation now always runs this arm; historical \texttt{jpeg\_unified.json}
still records a migration path that accepted Q70-only when it already covered the
grid~(\Cref{sec:jpeg}).

\section{Application: Any-Angle Circ-Harmonic Latent Watermark}
\label{sec:app-anyangle}

The section below instantiates a \emph{zero-bit}, keyed presence watermark on
Stable Diffusion initial noise.
Generation draws i.i.d.\ Gaussian \(z_0\in\mathbb{R}^{4\times64\times64}\), applies a
keyed mid-annulus circ-harmonic field on latent channel~3 only, and runs the full
denoising pipeline to an RGB image.
Detection inverts the VAE once and scores angular harmonic energy on the same
spatial ring; it does not search rotation, sector index, or residual angle.
Reported metrics (F1 on clean watermarked images, rotation pass rate under
expand-and-resize, false positives on unmarked SD images, visual quality) are
\emph{application} numbers on this instance.
They are not proofs of~(C), (S), (P), or~(R).
The carrier carries no user message bits; the decision is approximately one bit
per image (present vs.\ absent under the keyed template).
Rotation attacks rotate the \emph{pixel} image with \texttt{expand=True} and
resize back to the original canvas (bicubic rotation, Lanczos resize); this is
not pad-and-crop and is not inverted at detect time.

\subsection{Key material and trapdoor support}
\label{sec:anyangle-key}

\paragraph{Secret key.}
The implementation stores a PRC/LDPC key object; the any-angle codec uses only its
canonical PEM string \(\mathrm{PEM}(\sk)\) as a seed for deterministic pseudorandom
objects.
No LDPC encode, parity check, or syndrome appears in embed or detect for this
scheme (\emph{symmetric}: the same \(\sk\) embeds and verifies).

\paragraph{Keyed PRNG.}
For a UTF-8 tag~\(T\), define
\[
\mathrm{RNG}(T)=\mathrm{PRNG}\bigl(\lfloor \mathrm{SHA256}(T)\rfloor_{64}\bigr),
\]
where the subscript denotes the first 64 hexadecimal digits interpreted as an integer.

\paragraph{Harmonic trapdoor \((\Omega,\Omega^{\mathrm c})\).}
Let angular orders be indexed \(m\in\{1,\ldots,M\}\) with \(M=12\).
The trapdoor pool is the zero-based index set
\[
\mathcal{P}=\{m-1:\ m\in\{M_{\mathrm{lo}},\ldots,M_{\mathrm{hi}}\}\},
\qquad M_{\mathrm{lo}}=3,\ M_{\mathrm{hi}}=12,
\]
so \(|\mathcal{P}|=10\).
Draw a uniform permutation \(\pi\) of \(\mathcal{P}\) from
\(\mathrm{RNG}(\mathrm{PEM}(\sk)\,\|\,\text{\texttt{any-ntru}}\,\|\,3\,\|\,12\,\|\,6)\).
Split \(\pi\) into
\[
\Omega=\{\pi(1),\ldots,\pi(K)\},\qquad
\Omega^{\mathrm c}=\{\pi(K+1),\ldots,\pi(10)\},
\qquad K=6.
\]
Indices in \(\Omega\) are \emph{on-support}; indices in \(\Omega^{\mathrm c}\) are
\emph{control}.
Both sets are sorted ascending before use.
Amplitudes for embedding use \(a_m=\mathbf{1}[m-1\in\Omega]\) (unit on-support, zero
off-support within \(1..M\)).

\paragraph{Phases.}
For watermark channel \(c=3\), draw
\(\phi_m\sim\mathrm{Unif}[0,2\pi)\) i.i.d.\ from
\(\mathrm{RNG}(\mathrm{PEM}(\sk)\,\|\,\text{\texttt{any-phase}}\,\|\,c)\).

\subsection{Geometry: mid annulus on the latent grid}
\label{sec:anyangle-geom}

Work on the \(64\times64\) spatial slice of each latent channel with integer
radius \(\hat r=\mathrm{round}(\sqrt{(y-c_y)^2+(x-c_x)^2})\) about the grid center
\((c_y,c_x)=(31.5,31.5)\) in pixel coordinates.

\paragraph{Payload annulus.}
Pixels with \(\hat r\in\{R_{\mathrm{in}},\ldots,R_{\mathrm{out}}\}\),
\(R_{\mathrm{in}}=10\), \(R_{\mathrm{out}}=18\), on channel~3 form the payload.
For each integer radius, collect flat indices sorted by polar angle
\(\theta=\mathrm{atan2}(y-c_y,x-c_x)\); concatenate radii
\(10,\ldots,18\) into one ordered list
\(\{( \theta_j, z_j)\}_{j=1}^{N}\) with \(N=792\) in the locked configuration
(rings with fewer than 16 pixels are skipped at embed; all nine radii are active).

\paragraph{Protected sets.}
Let \(\mathcal{R}_{\mathrm{pay}}\) be all flat indices on the payload annulus
(channel~3) and \(\mathcal{R}_{\mathrm{in}}\) be indices with \(\hat r<R_{\mathrm{in}}\)
on all four channels.
The \emph{repair pool} is
\[
\mathcal{U}=\{1,\ldots,4HW\}\setminus
\bigl(\mathcal{R}_{\mathrm{pay}}\cup\mathcal{R}_{\mathrm{in}}\bigr).
\]
Sign flips for pseudorandomness never touch \(\mathcal{R}_{\mathrm{pay}}\) or
\(\mathcal{R}_{\mathrm{in}}\) (face/core latent is frozen).

\subsection{Carrier: magnitude-weighted circ-harmonic field on \(z_0\)}
\label{sec:anyangle-carrier}

\paragraph{Initial noise.}
Draw \(z_0\in\mathbb{R}^{4\times H\times W}\) i.i.d.\ \(\mathcal{N}(0,1)\) with
\(H=W=64\) from the generation seed.
Only \(z_0\) is modified before denoising; no step-wise latent edits.

\paragraph{Angular drive.}
On each payload radius ring with index set \(I_r\), angles \(\theta_i\) and host
values \(z_i\), form the normalized harmonic field
\[
h(\theta)=\sum_{m=1}^{M} a_{m-1}\cos(m\theta+\phi_m),\qquad
\widetilde h(\theta)=\mathrm{clip}\!\left(\frac{h(\theta)}{\mathrm{RMS}(h)},\,-3,\,3\right),
\]
with \(\mathrm{RMS}(h)=\sqrt{\mathbb{E}[h(\theta)^2]}+\varepsilon\) over pixels on
that ring.
The locked field strength is \(\alpha=0.62\).
Update
\[
z_i \leftarrow z_i + \alpha\,|z_i|\,\widetilde h(\theta_i),
\qquad i\in I_r,\ c=3.
\]
Other channels remain the original Gaussian draw except for subsequent repair on
\(\mathcal{U}\).

\paragraph{Pseudorandom sign repair (embed side only).}
Let \(\mathrm{flat}(z)\) be the rasterized latent.
Using \(\mathrm{RNG}(\mathrm{PEM}(\sk)\,\|\,\text{\texttt{:anyangle:uniform\_bits}})\),
apply up to six rounds of keyed single-index sign flips on indices in~\(\mathcal{U}\),
with occasional mean-rebalancing of \(\pm1\) counts on~\(\mathcal{U}\), until the
\emph{uniform-bits} score on~\(\mathcal{U}\) reaches \(0.99\)
(development gate aligned with the PRC threat model).
This is an embed-time constraint, not a detector early-reject rule.

\paragraph{Generative forward.}
Feed the marked \(z_0\) as \texttt{latents} to Stable Diffusion
(\(512\times512\) RGB, fixed step count in experiments, e.g.\ 15 DDIM steps) with
prompt~\(p\).
The output image is \(I=\mathrm{SD}(p,z_0)\).

\subsection{Detection: one-shot harmonic PSD and keyed contrast}
\label{sec:anyangle-detect}

\begin{algorithm}[t]
\caption{Blind detection (any-angle circ-harmonic NTRU)}
\label{alg:anyangle-detect}
\begin{algorithmic}[1]
\Require RGB image \(I\), secret key \(\sk\), VAE encoder \(\mathcal{E}\),
  threshold \(\tau\), layout \((R_{\mathrm{in}},R_{\mathrm{out}},c=3,M,K)\)
\State \(\hat z \leftarrow \mathcal{E}(I)\) \Comment{single VAE encode, no diffusion invert}
\State Recompute \((\Omega,\Omega^{\mathrm c})\) from \(\mathrm{PEM}(\sk)\) (\Cref{sec:anyangle-key})
\State Extract payload samples \(\{(\theta_j,\hat z_j)\}_{j=1}^{N}\) on channel~3,
  annulus \([R_{\mathrm{in}},R_{\mathrm{out}}]\)
\State Build \(z\in\mathbb{R}^N\) from \(\hat z_j\); optionally magnitude-weight:
  \(z_j\leftarrow \hat z_j\sqrt{|\hat z_j|+\varepsilon}\) then scale to unit mean weight
\State \(z\leftarrow (z-\bar z)/(\|z\|_2+\varepsilon)\)
\For{$m=1,\ldots,M$}
  \State \(c_m \leftarrow \frac{1}{N}\sum_j z_j e^{-im\theta_j}\) \Comment{discrete angular Fourier}
  \State \(\mathrm{PSD}[m-1] \leftarrow |c_m|^2\)
\EndFor
\State \(s_1 \leftarrow \dfrac{\mean_{i\in\Omega}\mathrm{PSD}[i]-\mean_{i\in\Omega^{\mathrm c}}\mathrm{PSD}[i]}
  {\mean_{i\in\Omega}\mathrm{PSD}[i]+\mean_{i\in\Omega^{\mathrm c}}\mathrm{PSD}[i]+\varepsilon}\)
  \Comment{support contrast}
\State Build template \(g_i=+1\) for \(i\in\Omega\), \(g_i=-1\) for \(i\in\Omega^{\mathrm c}\);
  \(s_2\leftarrow\) Pearson\((\mathrm{PSD}|_{\Omega\cup\Omega^{\mathrm c}}, g)\)
\State Repeat \(s_1,s_2\) with magnitude-weighted~\(z\); \(\mathrm{raw}\leftarrow\max\{s_1,s_2,\ldots\}\)
\State \(\mathrm{peak}\leftarrow \mathrm{clip}\bigl((\mathrm{raw}-0.15)/0.70,\,0,\,1\bigr)\)
\State \Return \(\mathbf{1}[\mathrm{peak}\ge\tau]\)
\end{algorithmic}
\end{algorithm}

\paragraph{Rotation invariance (ideal vs.\ deployed).}
For pure plane rotation about the latent center, each coefficient transforms as
\(c_m\mapsto e^{im\Delta\theta}c_m\), hence \(|c_m|\) and \(\mathrm{PSD}[m-1]\) are
exactly invariant.
Expand-and-resize on pixels, VAE nonlinearity, and denoising break this ideal;
the detector deliberately performs \emph{no} angle search, deskew, or center
synchronization (\(L=1\) extraction, one syndrome-free score).

\paragraph{Detection-call budget.}
Each verification uses exactly one VAE encode and one pass through
\Cref{alg:anyangle-detect} (four internal score branches collapsed by max).

\subsection{Rotation attack protocol (evaluation)}
\label{sec:anyangle-rot}

Given watermarked image \(I_w\), draw \(\Delta\sim\mathrm{Unif}[0,360)\) and form
\[
I_{\mathrm{rot}}=
\mathrm{Resize}_{512\times512}\!\Bigl(
\mathrm{Rotate}_{\mathrm{bicubic}}(I_w,\Delta;\,\texttt{expand=True})
\Bigr).
\]
Verify with \(\sk\) on \(I_{\mathrm{rot}}\) as in \Cref{alg:anyangle-detect} without
counter-rotation.
This matches the repository rotation-eval contract (attack-then-verify-as-is).

\subsection{Threshold calibration}
\label{sec:anyangle-calib}

Every evaluation run freezes a scalar threshold~\(\tau\) before scoring test images.

\paragraph{Default floor.}
\(\tau_{\mathrm{floor}}=0.5\) on the clipped peak scale of \Cref{alg:anyangle-detect}.

\paragraph{Negative-driven rule (large-sample script).}
On \(n_{\mathrm{neg}}\) unmarked SD images generated with the same pipeline and prompt
distribution, collect peaks \(\{p_i^{\mathrm{neg}}\}\).
Let \(q_{0.85}\) be their 85th percentile.
Set
\[
\tau=\max\bigl(\tau_{\mathrm{floor}},\ q_{0.85}+0.03\bigr).
\]
Positives are watermarked images at zero rotation for F1-style summaries; rotation
robustness is reported separately as the fraction of attacked positives with
\(\mathrm{peak}\ge\tau\).

\paragraph{Pseudorandom embed gate (reporting).}
Before benchmark runs, the scheme may report a \emph{pseudorandom metric}:
uniform-bits score on~\(\mathcal{U}\) after embed-only on a dummy Gaussian latent,
required to pass the same \(0.99\) gate used during repair.
Reproducing PR tables requires this check; it is independent of~\(\tau\).

\subsection{Locked operating point}
\label{sec:anyangle-params}

The reference implementation (\texttt{pk\_ring\_anyangle}) uses
\[
(R_{\mathrm{in}},R_{\mathrm{out}},c,M,M_{\mathrm{lo}},M_{\mathrm{hi}},K,\alpha)
=(10,18,3,12,3,12,6,0.62),
\]
payload size \(N=792\), zero user payload bits, symmetric key,
\(\tau_{\mathrm{floor}}=0.5\), and SD~\(512\times512\) with latent
\(4\times64\times64\).
Historical large-sample benchmarks under expand-and-resize rotation report
order-of-magnitude clean pass \(\approx85\%\), rotation pass \(\approx47\%\) on
96 attacked draws (48 seeds, two random angles each), and F1 \(\approx0.91\) when
positives are defined on clean watermarked images only (\Cref{sec:anyangle-calib}).

\section{Application: Circ-Harmonic Ring NTRU (B3)}
\label{sec:app-b3}

The section below instantiates a \emph{zero-bit}, keyed presence watermark on
Stable Diffusion initial noise (codec \texttt{ring\_ntru}, catalog id~\texttt{sd\_b3}).
Generation draws i.i.d.\ Gaussian $z_0\in\mathbb{R}^{4\times64\times64}$, embeds a
keyed sparse angular-harmonic pattern on a \emph{spatial} payload annulus (unique
ring pixels, angle-sorted), optionally repairs off-payload signs for pseudorandomness,
and runs the full denoising pipeline.
Detection performs one VAE encode and scores summed angular harmonic energy
$|c_m|^2$ on the same rings, with log-spectrum whitening and on-support vs.\
control contrast.
There is no rotation search, deskew, sector index, or LDPC syndrome at detect time
(the PRC key seeds trapdoor permutations only).
Reported metrics (F1, small-angle rotation under expand-and-resize, FPR on unmarked
SD images, visual quality) are application numbers on this instance, not proofs
of~(C), (S), (P), or~(R).
The carrier carries no user message bits; the decision is approximately one bit
per image (present vs.\ absent under the keyed template).

\subsection{Key material and NTRU-style support}
\label{sec:b3-key}

\paragraph{Secret key.}
As in~\Cref{sec:anyangle-key}, the stored PRC/LDPC key is used symmetrically;
B3 reads only $\mathrm{PEM}(\sk)$ for deterministic objects (no LDPC
encode/decode on this path).

\paragraph{Keyed PRNG.}
Define $\mathrm{RNG}(T)$ from $\mathrm{SHA256}(T)$ as in~\Cref{sec:anyangle-key}.

\paragraph{Harmonic orders.}
Let $M=24$ index angular orders $m\in\{1,\ldots,M\}$.
The NTRU trapdoor pool uses zero-based indices
\[
\mathcal{P}_{\mathrm N}=\{m-1:\ m\in\{M_{\mathrm{ntru}},\ldots,M\}\},
\qquad M_{\mathrm{ntru}}=10,
\]
so $|\mathcal{P}_{\mathrm N}|=14$.
Draw a uniform permutation $\pi$ of $\mathcal{P}_{\mathrm N}$ from
\[
\mathrm{RNG}\bigl(\mathrm{PEM}(\sk)\,\|\,\text{\texttt{ntru-rs}}\,\|\,24\,\|\,10\,\|\,\rho\bigr),
\qquad \rho=0.42.
\]
Set $K=\max\bigl(4,\lfloor \rho|\mathcal{P}_{\mathrm N}|\rfloor\bigr)$, capped at
$\lfloor|\mathcal{P}_{\mathrm N}|/2\rfloor$; in the locked configuration $K=6$.
Define on-support $\Omega=\{\pi(1),\ldots,\pi(K)\}$ and control
$\Omega^{\mathrm c}=\{\pi(K+1),\ldots,\pi(14)\}$ (both sorted).
Embedding amplitudes use $a_{m-1}=A_{\mathrm{on}}$ on $\Omega$ and
$a_{m-1}=A_{\mathrm{floor}}$ elsewhere, with $A_{\mathrm{on}}=3.35$,
$A_{\mathrm{floor}}=0$.

\paragraph{Phases and drive noise.}
On each watermark channel $c\in\mathcal{C}_{\mathrm{wm}}$, draw
$\phi_{m,c}\sim\mathrm{Unif}[0,2\pi)$ from
$\mathrm{RNG}(\mathrm{PEM}(\sk)\,\|\,\text{\texttt{ntru-phase}}\,\|\,c)$ and
i.i.d.\ $\xi\sim\mathcal{N}(0,\sigma_{\mathrm{drv}}^2)$ on ring pixels with
$\sigma_{\mathrm{drv}}=0.22$.

\subsection{Geometry: outer annulus on the latent grid}
\label{sec:b3-geom}

On each channel's $64\times64$ slice, integer radius
$\hat r=\mathrm{round}(\sqrt{(y-c_y)^2+(x-c_x)^2})$ about $(c_y,c_x)=(31.5,31.5)$.

\paragraph{Payload annulus.}
The default implementation uses
\[
R_{\mathrm{in}}=16,\qquad R_{\mathrm{out}}=28,\qquad
\mathcal{C}_{\mathrm{wm}}=\{0,1,2,3\}.
\]
For each $(c,\hat r)$ with $R_{\mathrm{in}}\le\hat r\le R_{\mathrm{out}}$, collect
flat indices sorted by $\theta=\mathrm{atan2}(y-c_y,x-c_x)$; concatenate over radii
and channels into payload index set $\mathcal{R}_{\mathrm{pay}}$ ($|\mathcal{R}_{\mathrm{pay}}|=7248$ in the four-channel lock).
Some benchmark scripts temporarily set $\mathcal{C}_{\mathrm{wm}}=\{3\}$ for
comparability with other carriers; the algorithm is unchanged.

\paragraph{Protected inner disk.}
$\mathcal{R}_{\mathrm{in}}=\{\text{pixels with }\hat r<R_{\mathrm{in}}\text{ on all channels}\}$.
Repair pool $\mathcal{U}=\{1,\ldots,4HW\}\setminus(\mathcal{R}_{\mathrm{pay}}\cup\mathcal{R}_{\mathrm{in}})$.

\subsection{Carrier: sign blend plus magnitude harmonic field}
\label{sec:b3-carrier}

\paragraph{Initial noise.}
Draw $z_0\in\mathbb{R}^{4\times H\times W}$, $H=W=64$, i.i.d.\ $\mathcal{N}(0,1)$;
only $z_0$ is modified before denoising.

\paragraph{Angular drive (per channel ring).}
On payload index set $I$ with angles $\{\theta_i\}$,
\[
h(\theta)=\sum_{m=1}^{M} a_{m-1}\cos(m\theta+\phi_{m,c})+\xi(\theta),\qquad
s_i=\mathrm{sign}\bigl(h(\theta_i)\bigr)\in\{\pm1\}\ \text{(0 $\mapsto$ $+1$)}.
\]

\paragraph{Sign blend.}
With blend $\beta=0.47$,
\[
z_i \leftarrow (1-\beta)\,z_i + \beta\,|z_i|\,s_i,\qquad i\in I.
\]

\paragraph{Magnitude field.}
Normalize $h$ without noise on ring $I$:
$\widetilde h=\mathrm{clip}(h/\mathrm{RMS}(h),-3,3)$.
Add
\[
z_i \leftarrow z_i + \eta\,|z_i|\,\widetilde h(\theta_i),\qquad \eta=0.26,\quad i\in I.
\]

\paragraph{Pseudorandom repair (embed only).}
Using $\mathrm{RNG}(\mathrm{PEM}(\sk)\,\|\,\text{\texttt{ring\_ntru:uniform\_bits}})$,
flip signs only at indices in~$\mathcal{U}$ (up to eight outer rounds, with
mean-rebalancing and lag-correlation reduction on~$\mathcal{U}$), until the
\emph{full} $4HW$ raster passes the uniform-bits gate with score $\ge 0.99$.
Separately, angular magnitude leakage on the payload ring is penalized so that
peak/mean of magnitude-weighted $|c_m|^2$ stays below a development ceiling
($\approx 3.5$ peak/mean ratio).

\paragraph{Generative forward.}
Marked $z_0$ is passed as \texttt{latents} to Stable Diffusion ($512\times512$ RGB,
fixed step count, e.g.\ 15) with prompt~$p$; output $I=\mathrm{SD}(p,z_0)$.

\subsection{Detection: whitened PSD and robust support contrast}
\label{sec:b3-detect}

\paragraph{Angular PSD (multi-channel).}
After $\hat z=\mathcal{E}(I)$, for each $c\in\mathcal{C}_{\mathrm{wm}}$ form samples
$\{(\theta_j,\hat z_{c,j})\}$ on that channel's annulus rings (same geometry as embed).
For each channel, compute discrete $|c_m|^2$ as in~\Cref{alg:anyangle-detect}
(optional sign-only or magnitude-weight preprocessing per branch), then \emph{sum}
PSD across channels into $\mathrm{PSD}\in\mathbb{R}^M$.

\paragraph{Log-spectrum whitening.}
Let $x_m=\log(\mathrm{PSD}[m-1]+\varepsilon)$. Convolve $\{x_m\}$ with a fixed
7-tap symmetric kernel (triangular weights) and divide
\[
\widetilde{\mathrm{PSD}}[m-1]
\leftarrow \frac{\mathrm{PSD}[m-1]}{\exp(\mathrm{smoothed\ log})+\varepsilon}.
\]

\paragraph{Support scores.}
With $(\Omega,\Omega^{\mathrm c})$ recomputed from $\mathrm{PEM}(\sk)$, define
mean contrast
\[
s_{\mathrm{mean}}=
\frac{\mean_{i\in\Omega}\widetilde{\mathrm{PSD}}[i]
-\mean_{i\in\Omega^{\mathrm c}}\widetilde{\mathrm{PSD}}[i]}
{\mean_{i\in\Omega}\widetilde{\mathrm{PSD}}[i]
+\mean_{i\in\Omega^{\mathrm c}}\widetilde{\mathrm{PSD}}[i]+\varepsilon},
\]
median contrast $s_{\mathrm{med}}$ (same formula with medians), and
$\mathrm{raw}_0=\max\{s_{\mathrm{mean}},s_{\mathrm{med}}\}$.

\paragraph{Branch fusion (NTRU only).}
Compute $\mathrm{raw}_0$ on float PSD; optionally replace with sign-thresholded PSD
if improvement exceeds $0.08$; optionally replace with magnitude-weighted float PSD
if improvement exceeds $0.20$. Let $\mathrm{raw}$ be the selected value.

\begin{algorithm}[t]
\caption{Blind detection (B3 / \texttt{ring\_ntru})}
\label{alg:b3-detect}
\begin{algorithmic}[1]
\Require RGB image $I$, secret key $\sk$, VAE encoder $\mathcal{E}$,
  threshold $\tau$, layout $(R_{\mathrm{in}},R_{\mathrm{out}},\mathcal{C}_{\mathrm{wm}},M,K)$
\State $\hat z \leftarrow \mathcal{E}(I)$
\State Recompute $(\Omega,\Omega^{\mathrm c})$ from $\mathrm{PEM}(\sk)$ (\Cref{sec:b3-key})
\State $\mathrm{PSD}\leftarrow 0$; for each $c\in\mathcal{C}_{\mathrm{wm}}$:
  accumulate $|c_m|^2$ on channel-$c$ annulus into $\mathrm{PSD}$
\State $\widetilde{\mathrm{PSD}}\leftarrow \mathrm{WhitenLog}(\mathrm{PSD})$
\State $\mathrm{raw}\leftarrow \mathrm{RobustContrast}(\widetilde{\mathrm{PSD}},\Omega,\Omega^{\mathrm c})$
  \Comment{mean/median on-support vs control}
\State Optionally refine $\mathrm{raw}$ using sign-PSD and mag-weight PSD branches
\State $\mathrm{peak}\leftarrow \mathrm{clip}\bigl((\mathrm{raw}+0.06)/0.68,\,0,\,1\bigr)$
\State \Return $\mathbf{1}[\mathrm{peak}\ge\tau]$
\end{algorithmic}
\end{algorithm}

\paragraph{Detection-call budget.}
One VAE encode and one fused score (\Cref{alg:b3-detect}); no angle or center search.

\subsection{Rotation attack protocol (evaluation)}
\label{sec:b3-rot}

Use the same expand-and-resize pixel rotation as~\Cref{sec:anyangle-rot} on
watermarked images.
Small-angle suites (e.g.\ $\pm1^\circ$--$\pm5^\circ$) report rotation pass rate
separately from clean-image F1; the detector never counter-rotates.

\subsection{Threshold calibration}
\label{sec:b3-calib}

\paragraph{Default floor.}
For \texttt{ring\_ntru} / \texttt{sd\_b3}, $\tau_{\mathrm{floor}}=0.48$ on the clipped
peak scale of \Cref{alg:b3-detect}.

\paragraph{Negative-driven rule (catalog / \texttt{calibrate\_threshold}).}
Collect peaks $\{p_i^{\mathrm{neg}}\}$ on $n_{\mathrm{neg}}$ unmarked SD images (or
Gaussian latents if no pipe).
Robustify: median/MAD inliers, $q_{0.85}$ on inliers, margins
($+0.018$ for NTRU), soft cap at $q_{0.95}+0.015$, optional anchor from
$n_{\mathrm{anchor}}=12$ unmarked Gaussian latents scored in latent space
($q_{0.95}+0.02$).
Set $\tau=\max(\tau_{\mathrm{floor}},\min\{\text{MAD/q/soft caps}\})$; ignore isolated
negative outliers above $\tau_{\mathrm{base}}+0.12$ when flagged.

\paragraph{Pseudorandom embed gate.}
Report $\min\{\text{uniform-bits score on full raster},\ \text{angular-mag score on payload}\}$ after embed on a dummy latent; require $\ge 0.99$ for PR tables.

\subsection{Locked operating point and relation to any-angle}
\label{sec:b3-params}

Reference \texttt{pk\_ring\_mass} (\texttt{ring\_ntru}) uses
\[
(R_{\mathrm{in}},R_{\mathrm{out}},\mathcal{C}_{\mathrm{wm}},M,M_{\mathrm{ntru}},\rho,
A_{\mathrm{on}},\beta,\eta)
=(16,28,\{0,1,2,3\},24,10,0.42,3.35,0.47,0.26),
\]
$K=6$ on-support harmonics among $m\ge 10$, payload size $7248$ (four channels),
zero user payload bits, symmetric key, $\tau_{\mathrm{floor}}=0.48$.
Compared to~\Cref{sec:app-anyangle}, any-angle moves the annulus inward
($10$--$18$), uses a single channel, shorter harmonic band ($M=12$, $m=3..12$),
\emph{sign-blend-free} magnitude-only embed, no log whitening, and a simpler
negative threshold rule---while sharing the same SD+$z_0$-only threat model and
one-shot $|c_m|^2$ philosophy.


\section{Experimental Protocol}
\label{sec:protocol}

The gates below are application criteria for the lattice watermark.
Passing them does not establish~\Cref{def:prcc}.

\paragraph{Data.}
COCO images are bilinearly resized to \(256\times256\).
Development splits (\texttt{tune}/\texttt{calib}/\texttt{test\_*}) are disjoint from
the hold-out confirmation split from COCO \texttt{train2017}
(\texttt{confirm\_splits.json}).

\paragraph{Attack.}
Unless noted: random rotation via pad-and-crop, then JPEG quality~70.
Residual-angle curves report mid-bin \(\delta\) points.
Additional translation after rotation+JPEG appears only in the boundary study.

\paragraph{Gates (test).}
Pass if (i)~random-\(\theta\) TPR \(\ge 90\%\),
(ii)~minimum mid-bin TPR \(\ge 90\%\), and
(iii)~Clopper--Pearson 95\% FPR upper bounds on unmarked and attacked negatives are
both \(\le 1\%\).
Calibration uses the stricter \(0.5\%\) empirical targets of~\Cref{sec:calib}.

\paragraph{Locked detector~C.}
\(m=64\), \(b=4\), \(w=1\), \(\eta=0\), \(L=8\), centers
\(\{(127.5,127.5),(128,128)\}\), lattice step \(s\approx 0.1416\) for the main
tables.
Each result table states its calibration rule and cites a JSON artifact under
\texttt{figures/}.

\section{Results}
\label{sec:experiments}

Numbers in this section are application metrics for the \(\eta=0\), \(w=1\)
watermark: distortion, rotation+JPEG true-positive rate, image false-positive upper
bounds, and detection calls.
They do not substitute for uniform-string soundness or pseudorandomness.

\subsection{Operating point and hold-out confirmation}
\label{sec:main}

\Cref{tab:psnr-scan} summarizes a distortion scan on a single development key
(\texttt{figures/psnr\_scan.json}; \textbf{calibration: Q70-only}).
Median PSNR \(\approx 38\,\mathrm{dB}\) (\(s\approx 0.1416\)) is the highest tested
setting that still passed all gates; \(39\,\mathrm{dB}\) and \(40\,\mathrm{dB}\) fail on
mid-bin TPR.

\begin{table}[t]
\centering
\caption{Distortion scan on detector~C (one development key; Q70-only calibration).
Source: \texttt{figures/psnr\_scan.json}.}
\label{tab:psnr-scan}
\begin{tabular}{@{}lcccc@{}}
\toprule
Target PSNR & Step \(s\) & Median PSNR & Mid-bin gate & Pass \\
\midrule
35 & 0.199 & 34.93 & pass & yes \\
36 & 0.178 & 35.93 & pass & yes \\
38 & 0.142 & 37.96 & pass & \textbf{yes} \\
39 & 0.128 & 38.84 & fail (\(0.888\)) & no \\
40 & 0.112 & 39.99 & fail (\(0.800\)) & no \\
\bottomrule
\end{tabular}
\end{table}

\Cref{tab:confirm} is the primary confirmation table: five keys, independent
per-key \textbf{Q70-only} calibration, confirmation-split images, attack $=$ random
rotation + JPEG~70 (\texttt{figures/confirm\_holdout.json}).
All five keys pass.
Per key: \(\approx 600\) marked positives for random-\(\theta\) TPR,
\(n=3000\) unmarked and \(n=3000\) attacked negatives for FPR.

\begin{table}[t]
\centering
\caption{Hold-out confirmation at \(\approx 38\,\mathrm{dB}\) (five keys; Q70-only
calibration).
FPR columns: Clopper--Pearson 95\% upper bounds.
Source: \texttt{figures/confirm\_holdout.json}.}
\label{tab:confirm}
\begin{tabular}{@{}rcccccc@{}}
\toprule
Key & \(\tau\) & PSNR & TPR & Mid min & FPR\textsubscript{none} & FPR\textsubscript{att} \\
\midrule
0 & 16 & 37.94 & 0.955 & 0.963 & 0.0026 & 0.0010 \\
1 & 17 & 37.96 & 0.965 & 0.975 & 0.0039 & 0.0044 \\
2 & 17 & 37.96 & 0.968 & 0.963 & 0.0030 & 0.0026 \\
3 & 17 & 37.97 & 0.967 & 0.950 & 0.0021 & 0.0026 \\
4 & 17 & 37.96 & 0.968 & 0.963 & 0.0030 & 0.0026 \\
\bottomrule
\end{tabular}
\end{table}

\subsection{Center-sync ablation}
\label{sec:center}

\Cref{tab:center-sync} isolates the half-pixel center effect
(\texttt{figures/center\_sync.json}; development step \(\approx 0.197\), three keys;
Q70-only).
Keeping the canvas center fixed (arm~A) or compensating with the true content center
(arm~B) removes the mid-bin valley.
Blind search over the two points
\(\{(127.5,127.5),(128,128)\}\) (arm~C) matches that recovery at \(2L=16\)
candidates.
Default-center extraction under pad-and-crop (arm~D) drops mid-bin TPR to about
\(0.82\).
Denser residual grids alone (\(L=32\)) do not close the valley
(\texttt{figures/dense\_search.json}).

\begin{table}[t]
\centering
\caption{Center-sync ablation (three keys, \(L=8\)).
Entries are minima across keys.
Source: \texttt{figures/center\_sync.json}.}
\label{tab:center-sync}
\begin{tabular}{@{}llcc@{}}
\toprule
Arm & Extract centers & Mid-bin TPR & Random-\(\theta\) TPR \\
\midrule
A & default (keep-center rot.) & 0.98 & 0.98 \\
B & oracle true center & 0.98 & 0.975 \\
C & \(\{(127.5,127.5),(128,128)\}\) & 0.98 & 0.975 \\
D & default under pad+crop & 0.82 & 0.915 \\
\bottomrule
\end{tabular}
\end{table}

\subsection{Unified JPEG threshold}
\label{sec:jpeg}

Per-quality recalibration shows operating points exist for each
\(Q\in\{50,\ldots,90\}\) (\texttt{figures/attack\_boundary.json}; per-\(Q\)
thresholds).

\paragraph{Historical unified-threshold artifact.}
\texttt{figures/jpeg\_unified.json} freezes one \(\tau\) per key and retests all
\(Q\).
In that recorded run, keys~0 and~1 used a Q70-calibrated threshold that already
covered \(Q\in\{50,\ldots,90\}\); key~2 failed Q70-only transfer at \(Q=60\) and was
recalibrated with mixed-Q, after which all three keys passed.
The repository protocol now \emph{always} executes mixed-Q for every key and retains
Q70-only as a migration arm; the JSON above has not been regenerated under that
always-mixed-Q regime.
We therefore treat the file as evidence that a single frozen \(\tau\) can cover the
tested quality grid (with mixed-Q needed for at least one key), not as a fresh
always-mixed-Q campaign.

\subsection{Fair comparison to PRC}
\label{sec:prc}

We match lattice carrier, two-center residual geometry, distortion, and images, and
replace only the code layer: block-circulant PRCC versus a sparse PRC detector that
ORs over \(m\) integer shifts at each of the \(2L\) geometry candidates
(\(2L\cdot m=1024\) code detections per image).
Both sides use \textbf{Q70-only} calibration after full search; PRC takes the lowest
legal score meeting the \(0.5\%\) caps; PRC detection is CSR-vectorized
(\texttt{figures/prc\_fair\_C.json}).

\begin{table}[t]
\centering
\caption{PRCC vs.\ PRC on matched carrier~C (three paired keys; means; Q70-only
calibration).
Source: \texttt{figures/prc\_fair\_C.json}.}
\label{tab:prc}
\begin{tabular}{@{}lccc@{}}
\toprule
Metric & PRCC & PRC & Ratio \\
\midrule
Code detections / image & 16 & 1024 & \(64\times\) \\
Mean code time (s) & 2.3 & 112.9 & \(\approx 50\times\) \\
Mean total time (s) & 102.7 & 215.5 & \(\approx 2.1\times\) \\
Mean TPR & 0.972 & 0.984 & --- \\
\bottomrule
\end{tabular}
\end{table}

The \(64\times\) call reduction is structural.
After vectorization, code-time speedup (\(\approx 50\times\)) tracks call count.
Total wall-clock speedup is only \(\approx 2.1\times\) because radial extraction
dominates.

\paragraph{False-positive footnote.}
PRCC key~2 records attacked negatives \(10/1500\) (95\% upper bound
\(\approx 1.13\%\)); we do not claim all three comparison keys meet a strict
\(1\%\) upper-bound budget.
Efficiency ratios are unaffected.

\subsection{Template + cyclic correlation baseline}
\label{sec:corr}

To place the method among classical geometric detectors, we keep the same carrier,
PSNR step, two-center \(\times L\) geometry, and images, and replace the code by a
secret \(\pm 1\) template with simultaneous-ring circular correlation (real FFT per
ring, sum lags; max over geometry).
Thresholds are calibrated after the full search (higher-is-better; lowest legal
score under the same \(0.5\%\) caps; Q70-only), matching the spirit of
search-and-correlate detectors~\cite{lin2001rst}
(\texttt{figures/corr\_fair\_C.json}).

\begin{table}[t]
\centering
\caption{PRCC vs.\ template+FFT correlation on matched carrier~C (three keys;
Q70-only; means).
A CORR detect call is one FFT cyclic correlation covering \(m\) lags.
Source: \texttt{figures/corr\_fair\_C.json}.}
\label{tab:corr}
\begin{tabular}{@{}lcc@{}}
\toprule
Metric & PRCC & Template+corr \\
\midrule
Detect calls / image & 16 & 16 \\
Mean TPR & 0.970 & 1.000 \\
Mean FPR\textsubscript{att} upper & 0.0047 & 0.0051 \\
Mean code time (s) & 2.1 & 2.1 \\
Mean total time (s) & 105.4 & 103.4 \\
\bottomrule
\end{tabular}
\end{table}

With FFT amortization, both detectors issue 16 score evaluations per image and have
essentially identical wall-clock cost.
On this matched carrier the template correlator achieves higher TPR (1.000 vs.\
0.970) at comparable FPR upper bounds.
Thus~\Cref{tab:prc} shows a structural call reduction against an \(m\)-fold PRC OR,
whereas~\Cref{tab:corr} shows that once integer lags are handled by FFT correlation,
the syndrome test does not improve detection rate or runtime under the present
protocol.
The distinctive measured benefit of the block-circulant layer remains the removal of
an explicit integer-sector loop relative to shift-OR code detectors---not a blanket
advantage over classical template correlation on the same geometry.

\subsection{Translation boundary}
\label{sec:shift}

After rotation+JPEG~70 we add a pixel translation
(\texttt{figures/attack\_boundary.json}; per-condition calibration).
Offsets \(\lvert\Delta\rvert\le 0.25\)~px often still admit a passing threshold; some
axis/key combinations slightly exceed the \(1\%\) FPR upper bound.
Offsets \(\lvert\Delta\rvert\ge 0.5\)~px yield no usable threshold under the locked
two-center detector.

\section{Discussion}
\label{sec:discussion}

\paragraph{What is established.}
At the coding layer, the block-circulant candidate meets noiseless integer-shift
completeness and, under full rank, the uniform-string tail in~\Cref{prop:S}.
At the application layer, the assembled detector---parity counts at \(w=1\), hard
radial QIM, and two-center residual search---meets the TPR/FPR/PSNR gates at
\(\approx 38\,\mathrm{dB}\) on five Q70-calibrated confirmation keys, and reduces
code-detection calls by \(64\times\) versus a shift-OR PRC baseline on the same
carrier.
A matched template+FFT correlator uses the same 16 geometry scores per image and
achieves higher TPR, so this watermark instance is not competitive with classical
correlation on raw detection rate.
That limit does not remove the \PRCC{} definition or the candidate's equivariance.

\paragraph{What is not established.}
The reported operating point uses \(\eta=0\) and therefore fails
pseudorandomness~(P) by the public parity constraint of~\Cref{sec:status}.
Image false-positive rates are not uniform-string soundness~(S), and the
experiments do not prove BSC robustness~(R).
Image-level undetectability is not claimed.
The mixed-Q formal protocol is implemented, but the archived unified-JPEG JSON still
reflects the older migration path; a full always-mixed-Q regenerate remains optional
bookkeeping.
General translation, scale, and generative inversion pipelines remain open.

\paragraph{Open directions.}
(1)~Noise \(\eta>0\) with a validated hardness story.
(2)~Translation-tolerant sync after oracle-center controls.
(3)~Richer classical baselines (e.g., closer reproductions of~\cite{lin2001rst}
log-polar projections) if claiming broad watermark competitiveness beyond the matched
carrier.
(4)~Distribution-preserving embedding in a diffusion pipeline~\cite{gunn2024undetectable}.

\section*{Acknowledgments}
Working draft. Numbers are traced to JSON artifacts under \texttt{figures/}.
Reference implementation: \texttt{experiments/src/}.

\appendix
\section{Failed full-circulant square matrix}
\label{app:failed}

\begin{construction}[Retracted full-circulant candidate]
Sample \(h\in\{0,1\}^n\) of weight \(w\) and let \(H\) be the \(n\times n\) circulant
with rows \(R_0(h),\ldots,R_{n-1}(h)\).
Encoding was described as sampling \(\ker H\) plus \(\Bern(\eta)\).
\end{construction}

\begin{proposition}[Trivial kernel at default parameters]
Let \(n=2^\ell\). Over \(\mathbb{F}_2\), \(X^n-1=(X+1)^n\).
If \(w\) is odd then \(h(1)=1\), so \(\gcd(h(X),X^n-1)=1\) and the circulant \(H\) is
invertible.
Hence \(\ker H=\{0\}\).
For \(n=256\), \(w=7\), encoding collapses to noise alone; small \(\eta\) distinguishes
codewords from uniform without the key.
\end{proposition}

\section{Early channel attempts}
\label{app:channel}

Painting i.i.d.\ signs on discrete FFT rings left half-bin grid-aligned bit-error
rates near \(1/2\) under pad-and-crop rotation.
Spatial annulus sectors with integer-only \(T_k\) rolls still failed at half-bin
residuals---motivating keyed residual search.
Margin (sign-push) carriers were less stable under JPEG than two-coset QIM.
Phase notes are archived under \texttt{notes/}.

\bibliographystyle{plain}
\bibliography{references}
\end{document}
